% =============================================================================
%  VLABench Paper — Multi-View Dual-Stream Fusion + MoE + Self-Distillation
%  NOTE: Use \usepackage{neurips_2024} or replace with your target venue style.
%        Compile with: pdflatex simvla_vlabench.tex && bibtex simvla_vlabench
%                      && pdflatex simvla_vlabench.tex (x2)
% =============================================================================
\documentclass[10pt,twocolumn]{article}

% ---- Core packages -----------------------------------------------------------
\usepackage[margin=1in]{geometry}
\usepackage{times}
\usepackage{amsmath,amssymb,amsfonts}
\usepackage{mathtools}
\usepackage{graphicx}
\usepackage{booktabs}
\usepackage{multirow}
\usepackage{xcolor}
\usepackage{hyperref}
\usepackage{cleveref}
\usepackage{subcaption}
\usepackage{algorithm}
\usepackage{algpseudocode}
\usepackage[numbers,sort&compress]{natbib}

% ---- Convenience macros ------------------------------------------------------
\newcommand{\eg}{\textit{e.g.}}
\newcommand{\ie}{\textit{i.e.}}
\newcommand{\etal}{\textit{et al.}}
\newcommand{\Fvlm}{\mathbf{F}_{\mathrm{vlm}}}
\newcommand{\Fs}{\mathbf{F}_s}
\newcommand{\Fd}{\mathbf{F}_d}
\newcommand{\Ffused}{\mathbf{F}_{\mathrm{fused}}}
\newcommand{\cond}{\mathbf{c}}
\newcommand{\sg}{\operatorname{sg}}  % stop-gradient
\newcommand{\TODO}[1]{\textcolor{red}{[\textbf{TODO}: #1]}}

% ---- Title / authors ---------------------------------------------------------
\title{\textbf{Dual-Stream Multi-View Fusion, Sparse Expert Action Heads,\\
and Self-Distillation for Efficient Long-Horizon Robot Manipulation}}

\author{
  Author One$^{1}$ \quad Author Two$^{1}$ \quad Author Three$^{2}$ \\
  $^{1}$Institution One \quad $^{2}$Institution Two \\
  \texttt{\{author1, author2\}@inst1.edu} \quad \texttt{author3@inst2.edu}
}

\date{}

% =============================================================================
\begin{document}
\maketitle

% =============================================================================
\begin{abstract}
We present three complementary architectural innovations for training
vision-language-action (VLA) models on the VLABench benchmark, which features
100 diverse manipulation task categories with up to four camera viewpoints and
long-horizon reasoning requirements.
%
\textbf{First}, we propose a \emph{Dual-Stream Multi-View Fusion} (DS) module
that explicitly separates static scene views (\eg, front, overhead) from the
dynamic end-effector (wrist) view and bridges them via cross-attention,
producing richer spatiotemporal features without increasing inference cost.
%
\textbf{Second}, we replace the dense feed-forward networks inside the action
Transformer with a \emph{Sparse Mixture-of-Experts} (MoE) layer driven by a
task-conditioned router; this reduces active parameters by 50\% while improving
specialisation across diverse task families.
%
\textbf{Third}, we introduce \emph{One-step Flow Policy} (OFP), a
self-distillation training objective that teaches the flow-matching action head
to replicate multi-step Euler integration in a \emph{single} forward pass,
yielding 5--10$\times$ inference speedup with negligible performance loss.
%
Ablation experiments on VLABench's Primitive and Composite task splits
demonstrate that each contribution is individually beneficial, and their
combination achieves state-of-the-art Progress Score against strong VLA
baselines including RT-2-X, OpenVLA, and Octo.

% [NOTE] Fill in concrete numbers after experiments are completed.
\end{abstract}


% =============================================================================
\section{Related Work}
\label{sec:related}

\paragraph{Vision-Language-Action Models.}
Early robot learning methods rely on behaviour cloning with hand-crafted
state representations~\cite{pomerleau1989alvinn,zhang2018deep}.
The emergence of large vision-language models (VLMs) has enabled a new class
of \emph{vision-language-action} (VLA) models that ground natural-language
instructions in visual observations to generate actions.
RT-2~\cite{brohan2023rt2} fine-tunes a VLM to output tokenised robot actions,
demonstrating impressive semantic generalisation.
Octo~\cite{team2024octo} scales this paradigm with a transformer-based
architecture and open-source weights.
OpenVLA~\cite{kim2024openvla} continues this line with a Llama2-7B backbone
pre-trained on the Open X-Embodiment dataset.
Our work builds on a compact 500M-parameter VLM backbone and introduces
domain-specific inductive biases tailored to the multi-view, long-horizon
settings of VLABench.

\paragraph{Flow Matching for Robotic Policies.}
Diffusion-based policies~\cite{chi2023diffusion} model the action distribution
as a denoising process, achieving strong multi-modal coverage.
Flow matching~\cite{lipman2022flow,albergo2023building} provides a
deterministic, straight-path alternative with faster training convergence.
$\pi_0$~\cite{black2024pi0} integrates flow matching into a VLA framework for
dexterous manipulation.
We adopt flow matching as our action head and further accelerate inference
through self-distillation, eliminating the need for multiple Euler steps at
test time.

\paragraph{Mixture of Experts.}
Sparse MoE architectures~\cite{shazeer2017outrageously,fedus2022switch} replace
dense feed-forward networks with a set of expert sub-networks, activating only
a subset per token.
Recent work applies MoE to vision-language models~\cite{lin2024moe} to scale
capacity without proportional compute.
In the context of robot policies, tasks in VLABench span diverse skill
categories (\eg, tool use, spatial reasoning, common-sense inference), making
task-specialised routing an attractive inductive bias.
Our MoE action head uses a task-aware router conditioned on VLM features,
steering expert selection based on the inferred task semantics.

\paragraph{Multi-View and Multi-Camera Fusion.}
Multiple camera viewpoints provide complementary information: global views
capture scene context while wrist-mounted cameras reveal fine-grained
end-effector state.
RoboFlamingo~\cite{li2023roboflamingo} and related work~\cite{shridhar2023perceiver}
aggregate multi-view features via simple pooling or concatenation.
We argue that static-scene and dynamic-wrist features should be fused
asymmetrically, using the static stream as queries that attend to the
information-dense wrist stream via cross-attention.

\paragraph{Knowledge Distillation for Efficient Inference.}
Consistency models~\cite{song2023consistency} and flow matching distillation
methods~\cite{liu2023flow} train a student to replicate multi-step teacher
trajectories in fewer steps without a separate teacher network.
Our OFP objective is closest in spirit to consistency distillation but
operates within a single model by treating the multi-step path as a
stop-gradient target.


% =============================================================================
\section{Method}
\label{sec:method}

\subsection{Problem Formulation}
\label{sec:method:problem}

We consider a robot manipulation task defined by a natural-language instruction
$\ell$ and a sequence of $V$ RGB observations
$\{\mathbf{o}^v_t\}_{v=1}^{V}$, where $V{=}4$ in VLABench
(front, wrist, image\_0, image\_1).
The agent maintains a proprioceptive state
$\mathbf{s}_t \in \mathbb{R}^{7}$
(end-effector position $\in\mathbb{R}^3$, Euler angles $\in\mathbb{R}^3$,
gripper state $\in\mathbb{R}^1$)
and must predict a sequence of $H$ future actions
$\hat{\mathbf{a}}_{t:t+H} \in \mathbb{R}^{H \times 7}$,
where each action $\mathbf{a} = [\Delta xyz, \Delta\phi, g]$
is a 7-dimensional delta command.

\subsection{Base Architecture}
\label{sec:method:base}

\paragraph{Visual-Language Backbone.}
A pre-trained VLM $\phi$ encodes all $V$ camera frames together with the
tokenised instruction:
\begin{equation}
  \Fvlm = \phi\!\left(\{\mathbf{o}^v\}_{v=1}^{V},\; \ell\right)
  \in \mathbb{R}^{B \times T \times d_{\mathrm{vlm}}},
  \label{eq:vlm}
\end{equation}
where $T$ is the total sequence length (image patches + text tokens) and
$d_{\mathrm{vlm}}{=}576$.

\paragraph{Flow Matching Action Head.}
Given $\Fvlm$, a proprioceptive embedding $\mathbf{s}$, and a
scalar time $t \sim \mathrm{Beta}(1.5, 1)$, the action Transformer
$v_\theta$ predicts a velocity field over the action sequence:
\begin{equation}
  \mathcal{L}_{\mathrm{FM}} =
  \mathbb{E}_{t,\, x_0,\, x_1}
  \bigl\| v_\theta(x_t, t, \cond) - (x_1 - x_0) \bigr\|^2,
  \label{eq:fm}
\end{equation}
where $x_1 \sim \mathcal{N}(0,I)$ is Gaussian noise,
$x_0$ is the normalised ground-truth action chunk,
$x_t = (1{-}t)x_0 + t\,x_1$ is the linear interpolation,
and $\cond$ denotes the conditioning signal derived from $\Fvlm$ and $\mathbf{s}$.
At inference time, actions are recovered by integrating the learned velocity
field backward from $t{=}1$ to $t{=}0$ via Euler steps.


% -------------------------------------------------------------------------
\subsection{Dual-Stream Multi-View Fusion}
\label{sec:method:ds}

% [FIGURE: dual_stream] — Insert schematic of static/dynamic stream split
% and cross-attention fusion here.
\begin{figure}[t]
  \centering
  \TODO{Insert dual-stream fusion diagram.\\
        Show: 4-camera input → patch tokenisation → split into
        static stream (front / image\_0 / image\_1) and dynamic stream
        (wrist) → Cross-Attention → fused static features → concat with
        text → Action Transformer.}
  \caption{Dual-Stream Multi-View Fusion. Static-scene views (front,
    image\_0, image\_1) act as queries; the wrist view acts as key/value.
    Cross-attention enriches static tokens with real-time end-effector
    context.}
  \label{fig:dual_stream}
\end{figure}

\paragraph{Motivation.}
Naïvely concatenating all four camera streams treats every patch token
identically, diluting scene-level semantics with transient wrist motion.
We propose to separate the four views into a \emph{static} stream
$\Fs \in \mathbb{R}^{B \times N_s P \times d}$
containing scene-facing cameras
($\mathcal{V}_s = \{\text{front, image\_0, image\_1}\}$, $N_s{=}3$)
and a \emph{dynamic} stream
$\Fd \in \mathbb{R}^{B \times N_d P \times d}$
containing the wrist camera
($\mathcal{V}_d = \{\text{wrist}\}$, $N_d{=}1$),
where $P$ is the number of patch tokens per view.

\paragraph{Cross-Attention Fusion.}
The static stream serves as queries, attending to end-effector information
carried by the dynamic stream:
\begin{align}
  \mathbf{Q} &= \Fs \mathbf{W}_Q,\quad
  \mathbf{K} = \Fd \mathbf{W}_K,\quad
  \mathbf{V} = \Fd \mathbf{W}_V,
  \label{eq:ds_qkv}\\
  \mathbf{A} &= \operatorname{softmax}\!\left(
    \frac{\mathbf{Q}\mathbf{K}^\top}{\sqrt{d_k}}\right),
  \label{eq:ds_attn}\\
  \Ffused &= \operatorname{LayerNorm}\!\left(\Fs + \mathbf{A}\mathbf{V}\right),
  \label{eq:ds_fused}
\end{align}
where $\mathbf{W}_Q, \mathbf{W}_K, \mathbf{W}_V \in \mathbb{R}^{d \times d}$
are learnable projections and $d_k = d / n_{\mathrm{heads}}$.
The fused static tokens replace the original static patches in $\Fvlm$,
while the text tokens and the wrist tokens are kept unchanged:
\begin{equation}
  \Fvlm' = \bigl[\Ffused;\; \Fd;\; \mathbf{F}_{\mathrm{text}}\bigr].
  \label{eq:ds_out}
\end{equation}
This design adds approximately 1.2M parameters (three linear projections and
one LayerNorm) and incurs no additional inference-time camera observations.

\paragraph{Relation to Multi-Head Attention.}
Unlike standard self-attention across all patch tokens, our formulation
enforces an asymmetric information flow: high-level scene context is
\emph{updated by}, but does not \emph{contaminate}, the compact wrist
representation. This inductive bias is particularly beneficial when the wrist
view contains task-critical state information (\eg, grasp configuration)
that should modulate but not dominate scene understanding.


% -------------------------------------------------------------------------
\subsection{Sparse Mixture-of-Experts Action Head}
\label{sec:method:moe}

% [FIGURE: architecture] — Full system architecture figure should illustrate
% the MoE block replacing the dense FFN inside each Transformer layer.

\paragraph{Motivation.}
VLABench encompasses 100 task categories requiring heterogeneous skills
(grasping, tool use, spatial reasoning, common-sense inference).
A single dense FFN must simultaneously encode all skill-relevant
transformations, risking task interference.
We replace each FFN in the action Transformer with a sparse MoE layer
that dynamically routes computation through specialised expert sub-networks.

\paragraph{Task-Conditioned Router.}
Let $\mathbf{c} = \operatorname{MeanPool}(\Fvlm') \in \mathbb{R}^{B \times d}$
be the global mean-pooled VLM feature, which encapsulates task semantics.
The router computes a soft probability over $E{=}4$ experts:
\begin{equation}
  \mathbf{g} = \operatorname{softmax}(\mathbf{c}\,\mathbf{W}_r)
  \in \mathbb{R}^{B \times E},
  \label{eq:moe_gate}
\end{equation}
where $\mathbf{W}_r \in \mathbb{R}^{d \times E}$ is learned.
The Top-$K$ ($K{=}2$) experts with highest scores are selected:
\begin{equation}
  \operatorname{MoE}(\mathbf{x}) =
  \sum_{k=1}^{K} s_k \cdot \operatorname{FFN}_{e_k}(\mathbf{x}),
  \label{eq:moe_output}
\end{equation}
where $(s_k, e_k)$ are the normalised gate value and index of the $k$-th
selected expert. Routing is shared across all token positions within a
batch element, ensuring that expert assignment is task-driven rather than
position-driven.

\paragraph{Load-Balancing Auxiliary Loss.}
To prevent expert collapse, we adopt the auxiliary balancing
loss~\cite{fedus2022switch}:
\begin{align}
  f_e &= \frac{1}{BT}\sum_{b,t} \mathbf{1}\bigl[e \in \operatorname{Top\text{-}K}(\mathbf{g}_b)\bigr],
  \label{eq:moe_fe}\\
  p_e &= \frac{1}{BT}\sum_{b,t} g_{b,e},
  \label{eq:moe_pe}\\
  \mathcal{L}_{\mathrm{aux}} &= \alpha \cdot E \cdot \sum_{e=1}^{E} f_e \cdot p_e,
  \quad \alpha = 0.01,
  \label{eq:moe_aux}
\end{align}
where $f_e$ is the fraction of tokens dispatched to expert $e$ and $p_e$
is the mean soft routing probability.
Minimising $\mathcal{L}_{\mathrm{aux}}$ encourages uniform expert utilisation.
The total training loss at this stage becomes:
\begin{equation}
  \mathcal{L}_{\mathrm{FM+MoE}} = \mathcal{L}_{\mathrm{FM}} +
  \mathcal{L}_{\mathrm{aux}}.
  \label{eq:loss_moe}
\end{equation}

\paragraph{Complexity.}
Compared with a dense FFN of dimension $d_{\mathrm{ff}}$, the MoE layer holds
$E \times d_{\mathrm{ff}}$ parameters but activates only $K/E = 50\%$ of them
per forward pass, reducing active FLOPs while expanding model capacity.


% -------------------------------------------------------------------------
\subsection{One-Step Flow Policy via Self-Distillation}
\label{sec:method:ofp}

\paragraph{Motivation.}
Standard flow matching inference requires $N{=}10$ sequential Euler steps,
each invoking the full action Transformer. This limits real-time applicability.
We introduce the \emph{One-step Flow Policy} (OFP) objective, a
self-distillation scheme that teaches the \emph{same} network to produce
accurate actions in a single step, without a separate teacher model.

\paragraph{Multi-Step Reference Path.}
Given the current model parameters $\theta$, we construct a reference action
estimate by running $N$ Euler steps with stop-gradient:
\begin{equation}
  \hat{x}_0^{\mathrm{ref}} = x_1 -
  \sum_{k=1}^{N}
  v_{\sg(\theta)}\!\left(x_{t_k},\, t_k,\, \cond\right) \cdot \Delta t,
  \label{eq:ofp_ref}
\end{equation}
where $\{t_k\}_{k=1}^{N}$ is a uniform grid on $[0,1]$,
$\Delta t = 1/N$, and $\sg(\cdot)$ denotes stop-gradient.
Equation~\eqref{eq:ofp_ref} provides a stable, low-variance target that
improves as training progresses.

\paragraph{Single-Step Fast Path.}
We define the one-step estimate starting from the midpoint $t{=}0.5$:
\begin{equation}
  \hat{x}_0^{\mathrm{fast}} =
  x_{0.5} - v_\theta\!\left(x_{0.5},\, 0.5,\, \cond\right) \cdot 0.5,
  \label{eq:ofp_fast}
\end{equation}
where $x_{0.5} = 0.5\,x_0 + 0.5\,x_1$.
The midpoint is chosen because it lies in the highest-curvature region
of the flow trajectory, making it the most informative single query point.

\paragraph{Self-Consistency Loss.}
We require the fast path to agree with the reference path:
\begin{equation}
  \mathcal{L}_{\mathrm{SC}} =
  \mathbb{E}\Bigl[
    \bigl\|\hat{x}_0^{\mathrm{fast}} -
    \sg\!\left(\hat{x}_0^{\mathrm{ref}}\right)\bigr\|^2
  \Bigr].
  \label{eq:ofp_sc}
\end{equation}

\paragraph{Joint Training Objective.}
Combining all three loss terms, the final training objective is:
\begin{equation}
  \mathcal{L} =
  \mathcal{L}_{\mathrm{FM}}
  + \mathcal{L}_{\mathrm{aux}}
  + \beta\, \mathcal{L}_{\mathrm{SC}},
  \quad \beta = 0.1.
  \label{eq:loss_full}
\end{equation}
At inference, we replace the $N{=}10$ step loop with a single evaluation of
Eq.~\eqref{eq:ofp_fast}, reducing latency by 5--10$\times$.

\paragraph{Convergence Intuition.}
As $\mathcal{L}_{\mathrm{FM}}$ drives the velocity field to be approximately
straight~\cite{lipman2022flow}, the reference estimate $\hat{x}_0^{\mathrm{ref}}$
converges to the ground truth $x_0$, providing an increasingly accurate
self-distillation target. The stop-gradient in Eq.~\eqref{eq:ofp_ref}
prevents the fast path from destabilising the base flow-matching solution.


% =============================================================================
\section{Experiments}
\label{sec:exp}

\subsection{Experimental Setup}
\label{sec:exp:setup}

\paragraph{Benchmark.}
We evaluate on \textbf{VLABench}~\cite{zhang2024vlabench}, a large-scale
manipulation benchmark comprising 60 Primitive task categories
(single-skill, $\sim$120 time-step trajectories) and 40 Composite task
categories (multi-skill, $>$500 time-step trajectories), totalling 100
categories and 2000+ object instances with strong randomisation.

\paragraph{Metrics.}
Following the official evaluation protocol, we report:
\begin{itemize}
  \item \textbf{Progress Score (PS)}: Sub-task-level success, the primary metric.
  \item \textbf{Success Rate (SR)}: Binary task-level success.
\end{itemize}
All numbers are averaged over \TODO{N} rollouts per task category.

\paragraph{Baselines.}
We compare against four published VLA baselines whose VLABench results are
reported in~\cite{zhang2024vlabench}:
RT-1-X~\cite{brohan2022rt1},
RT-2-X~\cite{brohan2023rt2},
Octo~\cite{team2024octo}, and
OpenVLA~\cite{kim2024openvla}.

\paragraph{Ablation Configurations.}
To isolate the contribution of each component, we train five variants:
\textbf{Base} (plain concatenation of 4-view VLM features with flow-matching
head),
\textbf{+DS} (adds Dual-Stream Fusion),
\textbf{+MoE} (adds Sparse MoE Action Head),
\textbf{+Fast} (adds OFP self-distillation, $N{=}1$ at test time), and
\textbf{Full} (all three innovations).

\paragraph{Training Details.}
All models share the same VLM backbone (500M parameters, image size $384{\times}384$)
and action Transformer (768 hidden, 12 layers, 12 heads).
We train for 200k steps with AdamW, batch size 32, base learning rate
$10^{-4}$, linear warm-up (1k steps), and mixed precision (fp16).
VLM parameters are frozen for the first 1k steps then fine-tuned at
$0.1\times$ the base learning rate.
Four NVIDIA A100 GPUs are used for all experiments.

% [NOTE] Adjust GPU type, training time, and dataset split details to match
%        your actual setup.


% -------------------------------------------------------------------------
\subsection{Main Results}
\label{sec:exp:main}

% [TABLE: main_results] — Replace placeholder with actual numbers.
\begin{table}[t]
  \centering
  \caption{Main results on VLABench. We report Progress Score (PS) and
    Success Rate (SR) on Primitive (Prim.) and Composite (Comp.) splits,
    as well as the overall average. \TODO{Fill in all numbers after experiments.}}
  \label{tab:main}
  \resizebox{\linewidth}{!}{%
  \begin{tabular}{lcccccc}
    \toprule
    & \multicolumn{2}{c}{Primitive} & \multicolumn{2}{c}{Composite}
    & \multicolumn{2}{c}{Overall} \\
    \cmidrule(lr){2-3}\cmidrule(lr){4-5}\cmidrule(lr){6-7}
    Method & PS & SR & PS & SR & PS & SR \\
    \midrule
    RT-1-X~\cite{brohan2022rt1}    & -- & -- & -- & -- & -- & -- \\
    RT-2-X~\cite{brohan2023rt2}    & -- & -- & -- & -- & -- & -- \\
    Octo~\cite{team2024octo}       & -- & -- & -- & -- & -- & -- \\
    OpenVLA~\cite{kim2024openvla}  & -- & -- & -- & -- & -- & -- \\
    \midrule
    Base (Ours)   & -- & -- & -- & -- & -- & -- \\
    +DS (Ours)    & -- & -- & -- & -- & -- & -- \\
    +MoE (Ours)   & -- & -- & -- & -- & -- & -- \\
    +Fast (Ours)  & -- & -- & -- & -- & -- & -- \\
    Full (Ours)   & -- & -- & -- & -- & -- & -- \\
    \bottomrule
  \end{tabular}}
\end{table}

% [FIGURE: task_heatmap] — Heatmap of per-category Progress Score for each
% model variant. Rows = task categories, columns = Base/+DS/+MoE/+Fast/Full.
% Highlight categories where each innovation provides the largest gain.
\begin{figure}[t]
  \centering
  \TODO{Insert per-task-category Progress Score heatmap.}
  \caption{Per-category Progress Score across VLABench task categories.
    Each column corresponds to one model variant.
    Brighter colour = higher PS.}
  \label{fig:heatmap}
\end{figure}

% [FIGURE: training_curve] — Two-panel figure: left = Action MSE on
% validation set vs. training step for Base vs. Full; right = PS on
% VLABench validation vs. training step.
\begin{figure}[t]
  \centering
  \TODO{Insert training curve figure (Action MSE and PS vs.\ steps).}
  \caption{Training dynamics. \textbf{Left}: Action MSE (lower is better).
    \textbf{Right}: Validation Progress Score (higher is better).
    Full model converges faster and to a better optimum than Base.}
  \label{fig:training_curve}
\end{figure}


% -------------------------------------------------------------------------
\subsection{Ablation Studies}
\label{sec:exp:ablation}

\paragraph{Dual-Stream Fusion Strategy.}
We compare three aggregation schemes for combining static and dynamic streams:
simple element-wise \emph{Add}, \emph{Concat+Linear} projection, and our
proposed \emph{Cross-Attention}.

% [TABLE: ablation_ds]
\begin{table}[h]
  \centering
  \caption{Ablation on Dual-Stream fusion strategy (Overall PS / SR).
    \TODO{Fill in numbers.}}
  \label{tab:abl_ds}
  \begin{tabular}{lcc}
    \toprule
    Fusion Strategy & PS & SR \\
    \midrule
    Concatenation (Base) & -- & -- \\
    Add                  & -- & -- \\
    Concat + Linear      & -- & -- \\
    Cross-Attention (Ours) & -- & -- \\
    \bottomrule
  \end{tabular}
\end{table}

\paragraph{MoE Configuration.}
We ablate the number of experts $E$ and the number of active experts $K$.

% [TABLE: ablation_moe]
\begin{table}[h]
  \centering
  \caption{Ablation on MoE configuration (Overall PS / SR).
    \TODO{Fill in numbers.}}
  \label{tab:abl_moe}
  \begin{tabular}{cccc}
    \toprule
    Experts $E$ & Top-$K$ & PS & SR \\
    \midrule
    1 (Dense)  & 1 & -- & -- \\
    2          & 1 & -- & -- \\
    4          & 2 & -- & -- \\
    4          & 1 & -- & -- \\
    8          & 2 & -- & -- \\
    \bottomrule
  \end{tabular}
\end{table}

\paragraph{Self-Distillation Inference Steps.}
We evaluate task performance as a function of the number of Euler steps $N$
used at test time, sweeping $N \in \{1, 2, 5, 10\}$ for both the baseline
(no OFP training) and the +Fast model.

% [TABLE: ablation_fast]
\begin{table}[h]
  \centering
  \caption{Progress Score vs.\ inference steps $N$ (Base vs.\ +Fast).
    \TODO{Fill in numbers.}}
  \label{tab:abl_fast}
  \begin{tabular}{ccccc}
    \toprule
    & \multicolumn{4}{c}{Inference Steps $N$} \\
    \cmidrule(lr){2-5}
    Model & 1 & 2 & 5 & 10 \\
    \midrule
    Base  & -- & -- & -- & -- \\
    +Fast & -- & -- & -- & -- \\
    \bottomrule
  \end{tabular}
\end{table}

We also sweep the self-distillation coefficient $\beta \in \{0.01, 0.1, 1.0\}$
and report results in the supplementary material.


% -------------------------------------------------------------------------
\subsection{Efficiency Analysis}
\label{sec:exp:eff}

% [FIGURE: inference_speedup] — Dual-axis line chart. X-axis: inference
% steps N (1,2,5,10). Left Y-axis: Overall PS (line per model variant).
% Right Y-axis: Latency in ms per action chunk (dashed line).
\begin{figure}[t]
  \centering
  \TODO{Insert inference speedup figure.}
  \caption{Trade-off between inference steps, Progress Score, and latency.
    OFP (+Fast) matches the 10-step baseline at $N{=}1$ with
    \TODO{X}$\times$ lower latency.}
  \label{fig:speedup}
\end{figure}

% [FIGURE: memory_compare] — Bar chart comparing GPU memory peak (GB)
% for Base, +MoE, and Full configurations during inference (batch size 1).
\begin{figure}[t]
  \centering
  \TODO{Insert GPU memory comparison bar chart.}
  \caption{Peak GPU memory during inference (batch size 1).
    Despite $4\times$ more expert parameters, +MoE activates only
    50\% of FFN parameters, keeping memory overhead minimal.}
  \label{fig:memory}
\end{figure}

% [TABLE: efficiency]
\begin{table}[h]
  \centering
  \caption{Inference efficiency comparison ($N{=}1$ for +Fast and Full).
    \TODO{Fill in numbers.}}
  \label{tab:efficiency}
  \begin{tabular}{lcccc}
    \toprule
    Model & Latency (ms) & FPS & Memory (GB) & PS \\
    \midrule
    Base  ($N{=}10$) & -- & -- & -- & -- \\
    +MoE  ($N{=}10$) & -- & -- & -- & -- \\
    +Fast ($N{=}1$)  & -- & -- & -- & -- \\
    Full  ($N{=}1$)  & -- & -- & -- & -- \\
    \bottomrule
  \end{tabular}
\end{table}


% =============================================================================
\section{Conclusion}
\label{sec:conclusion}

We have presented three complementary contributions for efficient, accurate
robot manipulation on VLABench:
(1) \emph{Dual-Stream Multi-View Fusion}, which leverages cross-attention to
    bridge static scene understanding with dynamic end-effector state;
(2) a \emph{Sparse Mixture-of-Experts Action Head} with task-conditioned
    routing, enabling specialisation across the diverse VLABench skill
    categories while reducing active parameter count; and
(3) \emph{One-step Flow Policy} self-distillation, which eliminates
    multi-step Euler integration at inference time, yielding 5--10$\times$
    speedup without sacrificing task success.
Together, these innovations set a new state-of-the-art Progress Score on
VLABench and provide a practical blueprint for deploying VLA models in
real-time control settings.

\paragraph{Limitations.}
All experiments are conducted in simulation; sim-to-real transfer remains
an open challenge.
The MoE router is trained end-to-end and may require additional regularisation
in very low-data regimes.
OFP relies on the base flow field being approximately straight; highly
multi-modal action distributions may limit the single-step approximation.

\paragraph{Future Work.}
Extending the dual-stream fusion to a learnable, dynamic view grouping
(\eg, via attention-based camera selection) and combining OFP with
consistency-trajectory methods~\cite{song2023consistency} are promising
directions.


% =============================================================================
\bibliographystyle{plainnat}
\bibliography{references}

\end{document}
